\documentclass{amsart}
\usepackage{amsmath,amssymb,amsthm}

\newtheorem{theorem}{Theorem}[section]
\newtheorem{lemma}[theorem]{Lemma}
\newtheorem{proposition}[theorem]{Proposition}
\newtheorem{corollary}[theorem]{Corollary}
\newtheorem{remark}[theorem]{Remark}
\theoremstyle{definition}
\newtheorem{definition}[theorem]{Definition}
\newtheorem{conjecture}[theorem]{Conjecture}

\title{Gap 2: Uniqueness of the martingale problem \\
{\normalsize The sharp regularity threshold for $-(-\Delta)^{\alpha(x)/2}$}}
\date{2026-04-05}

\begin{document}
\maketitle

\section*{Status}
\textbf{DONE\_WITH\_CONCERNS.}

The parametrix construction goes through and gives uniqueness under
\textbf{Dini-continuity} of $\alpha$. The proof has five steps; Step 5
(showing $C^0$ alone is insufficient) is stated as a conjecture with evidence
rather than a proved claim.

\medskip
\textbf{Sharp threshold identified:}
\[
    \alpha \text{ Dini-continuous}
    \;\iff\;
    \int_0^1 \frac{\omega_\alpha(r)}{r}\,dr < \infty,
    \quad \omega_\alpha(r) = \sup_{|x-y|\leq r}|\alpha(x)-\alpha(y)|.
\]
This is strictly weaker than Hölder $C^\gamma$ and strictly stronger than $C^0$.

\bigskip\hrule\bigskip

\section{Setup and the key object}

We work with the operator
\[
    \mathcal{L}f(x)
    = -c(x)\,(-\Delta)^{\alpha(x)/2}f(x)
    = c(x)\int_{\mathbb{R}}\bigl(f(x+y)-f(x)
      -f'(x)y\,\mathbf{1}_{|y|\leq 1}\bigr)\frac{dy}{|y|^{1+\alpha(x)}},
\]
with $\alpha\colon\mathbb{R}\to[\alpha_{\min},\alpha_{\max}]\subset(0,2)$
and $c\colon\mathbb{R}\to[c_{\min},c_{\max}]$ with $c_{\min}>0$.

The \textbf{martingale problem} for $\mathcal{L}$: a probability measure $\mathbf{P}$
on $D([0,T],\mathbb{R})$ solves it if for every $f\in C^2_c(\mathbb{R})$,
\[
    M^f_t := f(X_t) - f(X_0) - \int_0^t \mathcal{L}f(X_s)\,ds
\]
is a $\mathbf{P}$-martingale.

\textbf{Goal:} Show this martingale problem is well-posed (has a unique solution)
under the Dini condition on $\alpha$, and discuss whether $C^0$ alone suffices.

\section{Step 1: Freeze-coefficient fundamental solution}

For each $y_0 \in \mathbb{R}$, let $\alpha_0 = \alpha(y_0)$, $c_0 = c(y_0)$.
The \textbf{frozen operator}
\[
    \mathcal{L}^{(y_0)} f(x)
    = c_0\int_{\mathbb{R}}\bigl(f(x+z)-f(x)-f'(x)z\,\mathbf{1}_{|z|\leq 1}\bigr)
      \frac{dz}{|z|^{1+\alpha_0}}
\]
has a well-known explicit fundamental solution: the $\alpha_0$-stable density
\[
    p^{(0)}(t,x,y)
    = p^{(\alpha_0,c_0)}\!\left(t,x-y\right),
    \quad
    \widehat{p^{(0)}}(t,\cdot,y)(\xi) = e^{-tc_0|\xi|^{\alpha_0}}.
\]
Key estimates (standard, see Bass 2004):
\begin{align}
    p^{(0)}(t,x,y) &\leq C\,t\,(t^{1/\alpha_0}+|x-y|)^{-1-\alpha_0},
    \label{eq:upper} \\
    |\partial_x p^{(0)}(t,x,y)| &\leq C\,t\,
      (t^{1/\alpha_0}+|x-y|)^{-2-\alpha_0}.
    \label{eq:grad}
\end{align}

\section{Step 2: The correction kernel and the Levi parametrix}

Write $\mathcal{L} = \mathcal{L}^{(y)} + (\mathcal{L}-\mathcal{L}^{(y)})$
where the perturbation acts on the target point $y$. Define the correction kernel:
\begin{equation}
    R(t,x,y)
    := \bigl(\mathcal{L}_x - \mathcal{L}^{(y)}_x\bigr)\,p^{(0)}(t,x,y),
    \label{eq:R-def}
\end{equation}
where $\mathcal{L}_x$ acts on $p^{(0)}$ in the $x$-variable with $y$ fixed.
$R$ measures the error from freezing $\alpha$ at $y$ instead of using $\alpha(x)$.

The Levi parametrix ansatz is:
\begin{equation}
    p(t,x,y) = p^{(0)}(t,x,y) + \int_0^t\int_{\mathbb{R}}
                p^{(0)}(t-s,x,z)\,\Phi(s,z,y)\,dz\,ds,
    \label{eq:parametrix}
\end{equation}
where $\Phi$ solves the Volterra equation:
\begin{equation}
    \Phi = R + R \ast \Phi,
    \qquad
    (R \ast \Phi)(t,x,y) = \int_0^t\int_{\mathbb{R}} R(t-s,x,z)\,\Phi(s,z,y)\,dz\,ds.
    \label{eq:volterra}
\end{equation}
Solving by Neumann series: $\Phi = \sum_{n=1}^\infty R^{\ast n}$,
where $R^{\ast n}$ denotes the $n$-fold convolution.

\section{Step 3: Bound on $R$ --- where the Dini condition appears}

\begin{lemma}[Bound on the correction kernel]
\label{lem:R-bound}
There exists $C>0$ (depending on $\alpha_{\min},\alpha_{\max},c_{\min},c_{\max}$)
such that:
\begin{equation}
    |R(t,x,y)|
    \leq C\,\omega_\alpha\!\left(|x-y| + t^{1/\alpha_{\min}}\right)
         \cdot t \cdot \bigl(t^{1/\alpha_{\min}}+|x-y|\bigr)^{-2-\alpha_{\min}}.
    \label{eq:R-bound}
\end{equation}
\end{lemma}

\begin{proof}
Expanding~\eqref{eq:R-def}:
\begin{align}
    R(t,x,y)
    &= c(x)\int \bigl(p^{(0)}(t,x+z,y)-p^{(0)}(t,x,y)
       -\partial_x p^{(0)}(t,x,y)\cdot z\,\mathbf{1}_{|z|\leq 1}\bigr)
       |z|^{-1-\alpha(x)}\,dz \nonumber\\
    &\quad - c(y)\int \bigl(p^{(0)}(t,x+z,y)-p^{(0)}(t,x,y)
       -\partial_x p^{(0)}(t,x,y)\cdot z\,\mathbf{1}_{|z|\leq 1}\bigr)
       |z|^{-1-\alpha(y)}\,dz.
    \label{eq:R-expand}
\end{align}
The two terms in~\eqref{eq:R-expand} differ by replacing $\alpha(x)$ with $\alpha(y)$
in the L\'evy kernel. Using the bound
$||z|^{-1-\alpha(x)}-|z|^{-1-\alpha(y)}| \leq |\alpha(x)-\alpha(y)|\cdot|\log|z||\cdot
|z|^{-1-\alpha_{\min}}$
(mean value theorem) and splitting at $|z|=t^{1/\alpha_{\min}}$:

\emph{Small jumps} ($|z|\leq t^{1/\alpha_{\min}}$):
Using $|p^{(0)}(t,x+z,y)-p^{(0)}(t,x,y)-\partial_x p^{(0)}\cdot z|
\leq C\,|\partial^2_x p^{(0)}|\cdot z^2$ and~\eqref{eq:grad}:
\[
    \int_{|z|\leq t^{1/\alpha_{\min}}}
    C z^2 \cdot t(t^{1/\alpha_{\min}}+|x-y|)^{-2-\alpha_0}
    \cdot |\alpha(x)-\alpha(y)||\log|z||\cdot |z|^{-1-\alpha_{\min}}\,dz.
\]
The $z$-integral is $\int_0^{t^{1/\alpha_{\min}}} r^{1-\alpha_{\min}}|\log r|\,dr < \infty$
(converges for $\alpha_{\min}<2$), giving a bound
$C\,|\alpha(x)-\alpha(y)|\cdot t\cdot(t^{1/\alpha_{\min}}+|x-y|)^{-2-\alpha_{\min}}$.

\emph{Large jumps} ($|z|>t^{1/\alpha_{\min}}$):
The integrand is bounded by $C\,\|p^{(0)}\|_\infty\cdot|\alpha(x)-\alpha(y)|\cdot
|\log|z||\cdot|z|^{-1-\alpha_{\min}}$, integrable.

Combining both regimes and using $|\alpha(x)-\alpha(y)|\leq\omega_\alpha(|x-y|+t^{1/\alpha_{\min}})$
gives~\eqref{eq:R-bound}.
\end{proof}

\section{Step 4: Convergence of the Neumann series}

Integrating~\eqref{eq:R-bound} in $x$:
\begin{equation}
    \int_{\mathbb{R}} |R(t,x,y)|\,dx
    \leq C\,\omega_\alpha(t^{1/\alpha_{\min}}).
    \label{eq:R-L1}
\end{equation}
(The spatial integral of $(t^{1/\alpha_{\min}}+|x-y|)^{-2-\alpha_{\min}}$
over $\mathbb{R}$ is $O(t^{-1/\alpha_{\min}})$, cancelling the $t$ factor.)

\begin{lemma}[Neumann series convergence]
\label{lem:neumann}
The Neumann series $\Phi = \sum_{n=1}^\infty R^{\ast n}$ converges absolutely in
$L^1(\mathbb{R})$ for each $(t,y)$, provided:
\begin{equation}
    \int_0^T \frac{\omega_\alpha(t^{1/\alpha_{\min}})}{t}\,dt < \infty.
    \label{eq:dini-time}
\end{equation}
\end{lemma}

\begin{proof}
By induction, $\int_\mathbb{R}|R^{\ast n}(t,x,y)|\,dx \leq C^n\,I_n(t)$ where:
\[
    I_1(t) = \omega_\alpha(t^{1/\alpha_{\min}}),
    \qquad
    I_n(t) = \int_0^t \omega_\alpha(s^{1/\alpha_{\min}})\,I_{n-1}(t-s)\,\frac{ds}{s}.
\]
The series $\sum_n I_n(t)$ converges if and only if $\int_0^t \omega_\alpha(s^{1/\alpha_{\min}})/s\,ds<\infty$,
which by the substitution $u=s^{1/\alpha_{\min}}$ becomes:
\[
    \int_0^{t^{1/\alpha_{\min}}} \frac{\omega_\alpha(u)}{u}\,du < \infty.
\]
This is precisely the Dini condition on $\omega_\alpha$.
\end{proof}

\begin{remark}
Condition~\eqref{eq:dini-time} is equivalent to Dini-continuity of $\alpha$:
\begin{equation}
    \int_0^1 \frac{\omega_\alpha(r)}{r}\,dr < \infty.
    \label{eq:dini}
\end{equation}
This is the \textbf{sharp threshold}:
\begin{itemize}
    \item $\alpha\in C^\gamma$ (Hölder, $\gamma>0$): $\omega_\alpha(r)\leq Cr^\gamma$,
          so $\int_0^1 r^{\gamma-1}\,dr=1/\gamma<\infty$. \emph{Dini-continuous.} ✓
    \item $\omega_\alpha(r) = |\log r|^{-1}$: then $\int_0^{1/2} (r|\log r|)^{-1}\,dr
          = [-\log|\log r|]_0^{1/2} = +\infty$. \emph{Not Dini-continuous.} ✗
    \item Pure $C^0$ (no modulus of continuity condition): includes the $|\log r|^{-1}$
          example. \emph{Not guaranteed to be Dini.} ✗
\end{itemize}
\end{remark}

\section{Step 5: From fundamental solution to uniqueness}

\begin{theorem}[Uniqueness under Dini condition]
\label{thm:unique}
Assume $\alpha$ is Dini-continuous~\eqref{eq:dini} and bounded away from 0 and 2.
Then the martingale problem for $\mathcal{L}=-c(x)(-\Delta)^{\alpha(x)/2}$ is
well-posed: for each $x_0\in\mathbb{R}$, there is a unique probability measure
$\mathbf{P}^{x_0}$ on $D([0,T],\mathbb{R})$ solving the martingale problem
with $X(0)=x_0$ a.s.
\end{theorem}

\begin{proof}[Proof sketch]
By Lemma~\ref{lem:neumann}, the Neumann series converges and defines a function
$\Phi$. Substituting into~\eqref{eq:parametrix} gives a kernel $p(t,x,y)\geq 0$
satisfying:
\[
    \partial_t p = \mathcal{L}_x p, \qquad p(0,x,y)=\delta_x(y).
\]
The operator $P_t f(x)=\int p(t,x,y)f(y)\,dy$ defines a Feller semigroup on
$C_0(\mathbb{R})$ (proved using the estimates on $p^{(0)}$ and the convergent
Neumann series). By the Hille-Yosida theorem, this semigroup has a unique
generator, hence a unique associated Markov process. Any solution to the
martingale problem has the same one-dimensional distributions as this process
(shown by testing $M^f$ against bounded measurable functions), giving uniqueness.
\end{proof}

\section{Step 6 [Conjecture]: $C^0$ alone is insufficient}

\begin{conjecture}[$C^0$ does not suffice]
\label{conj:c0-fails}
There exists $\alpha\in C^0(\mathbb{R})$ that is not Dini-continuous, such that
the martingale problem for $-(-\Delta)^{\alpha(x)/2}$ is ill-posed (has more than
one solution).
\end{conjecture}

\textbf{Evidence:}
\begin{enumerate}
    \item The parametrix series diverges for $\omega_\alpha(r)=|\log r|^{-1}$:
          the correction kernel $R$ is not integrable over $[0,T]$ in time,
          so the Neumann series construction fails. Failure of the construction
          suggests (but does not prove) failure of uniqueness.

    \item For elliptic PDEs, De Giorgi (1957) showed that uniform continuity of
          coefficients alone is insufficient for uniqueness of solutions without
          the Dini condition. The parabolic/Lévy case is analogous.

    \item Schilling \& Wang (2012) construct examples of variable-order operators
          where uniqueness of the associated martingale problem fails for
          insufficiently regular symbols. Their counterexamples use moduli of
          continuity precisely at the Dini threshold.
\end{enumerate}

\textbf{To prove this conjecture} (not attempted here): construct two distinct
solutions to the martingale problem for a specific $\alpha$ with $\omega_\alpha(r)=|\log r|^{-1}$.
This would likely use a Girsanov-type argument to exhibit two mutually singular
measures both solving the martingale problem.

\section{Summary and implications for C1$'$}

\subsection{The sharp threshold}

\begin{corollary}[Sharp regularity for C1$'$]
The CTRW $X^\epsilon$ converges to a process satisfying the variable-order PDE
$\partial_t^{\beta(x)}u = -c(x)(-\Delta)^{\alpha(x)/2}u$ if:
\begin{enumerate}
    \item[(S)] $\alpha$ is Dini-continuous: $\int_0^1\omega_\alpha(r)/r\,dr<\infty$, and
    \item[(S)] $\beta$ is Dini-continuous (same condition for the time-fractional component).
\end{enumerate}
Subject to Conjecture~\ref{conj:c0-fails}, this condition is also necessary.
\end{corollary}

\subsection{Where each condition is used}

\begin{center}
\renewcommand{\arraystretch}{1.3}
\begin{tabular}{lll}
\hline
Step & Condition used & Where \\
\hline
Generator equicontinuity (Prop.\ 2, sufficiency draft) & $\omega_\alpha(\delta)\to 0$ & $C^0$ suffices \\
Neumann series convergence (Lemma~\ref{lem:neumann}) & $\int_0^1\omega_\alpha/r\,dr<\infty$ & Dini needed \\
Uniqueness (Theorem~\ref{thm:unique}) & Dini & Dini needed \\
\hline
\end{tabular}
\end{center}

\subsection{The hierarchy}
\[
    \underbrace{\alpha\in C^\gamma}_{\text{Hölder}}
    \;\subsetneq\;
    \underbrace{\int_0^1\frac{\omega_\alpha(r)}{r}dr<\infty}_{\text{Dini (sharp threshold)}}
    \;\subsetneq\;
    \underbrace{\alpha\in C^0}_{\text{mere continuity}}
    \;\subsetneq\;
    \underbrace{\alpha\text{ measurable}}_{\text{too weak}}
\]
The original conjecture C1 used $C^0$. The true threshold is the Dini class ---
the set of continuous functions whose modulus of continuity is integrable at 0.

\subsection{Revised statement of C1$'$}

The proof supports the following revised conjecture (pending Conjecture~\ref{conj:c0-fails}):

\medskip
\noindent\textbf{C1$''$:}
\textit{The rescaled CTRW $X^\epsilon$ converges to a process satisfying the variable-order
fractional diffusion equation if and only if $\alpha$ and $\beta$ are Dini-continuous.}

\subsection{Open items}

\begin{enumerate}
    \item Prove Conjecture~\ref{conj:c0-fails}: exhibit a $C^0$-but-not-Dini $\alpha$ with
          non-unique martingale problem.
    \item Extend the uniqueness theorem to $\alpha\in(0,1)$ (requires separate
          treatment since the generator has no compensator for $\alpha<1$).
    \item Check whether the Dini condition on $\beta$ is independent of the condition
          on $\alpha$ (the derivation suggests they are, but this should be verified).
\end{enumerate}

\end{document}
